\chapter{Metodologia}
\label{chap:metodologia}

\section{Etapas de Desenvolvimento do Modelo}
\label{sec:desenhopesquisa}
Este presente capítulo descreverá as etapas do desenvolvimento da Simulação Híbrida. De forma geral, a abordagem de eventos discretos (DE) simula as etapas macro gerais do processo no modelo computacional, enquanto a modelagem baseada em agentes (ABM) simula as decisões dos trabalhadores em relação à conclusão de cada atividade necessária. A simulação DE foi escolhida por ser considerada adequada para modelar o fluxo sequencial do canteiro de obras, considerando as mudanças de estados e recursos à medida que a sequência de eventos progride. Por outro lado, a ABM foi considerada satisfatória na representação do processo de tomada de decisão e dos padrões de comportamento dos trabalhadores inseridos no contexto dos eventos do canteiro de obras.


Considerando que as duas abordagens de simulação funcionam juntas neste estudo dentro da mesma plataforma de desenvolvimento, ambas seguem as etapas descritas na Figura 1, conforme originalmente definido por Crooks et al. (2019) para modelos ABM. De acordo com a Figura 1, o modelo deve ser preparado, projetado, implementado e submetido a procedimentos de verificação, calibração e validação antes das etapas de explicação e previsão. A Tabela 1, também baseada em Crooks et al. (2019), descreve o protocolo geral e os parâmetros definidos para o projeto do modelo. Este estudo concentra-se na modelagem do problema, e não em todos os elementos do sistema construtivo. Os dados foram coletados no Canteiro de Obras A.



\section{SED}
\label{sec:integraçao}